\section{Data provenance, certification, and the one open endpoint}
\label{app:data}

\paragraph{Sources.} All three studies use real Hyperliquid perpetual-futures observations.
Q17 draws on native on-change best-bid/offer contracts. Q18 draws on sampled-L2 receipt-clock
snapshots. The Q16 reconstruction diagnostic draws on certified BTC breakout and rebound
cohorts. Every experiment is gated by a source-and-schema audit that certifies event ordering
and receipt timing before the experiment is admitted, so a comparison runs only on inputs
whose validity has been checked. Raw market dumps are never redistributed. Each released
package ships compact metrics, checksums, and acquisition instructions instead, and no
committed path assumes access to the original development machine.

\paragraph{Restricted-data policy.} Raw and event-level derivatives are held under the data
provider's internal-analysis terms and are excluded from the release. Each package that
depends on such data states so and ships a license-clean smoke fixture so that its entry
point runs for anyone. Provider receipt time is uncalibrated against exchange time, and the
sampled-L2 feed does not establish native queue priority. These are recorded as limitations
next to every claim that could be affected.

\paragraph{The open native-FIFO endpoint.} The strictest form of Q16, native-FIFO count
tightening on a fully admitted order-lifecycle cohort, is not yet computed, and we state
this rather than approximate it. We located the exact public source: an open December 2025
Level-4 order-book archive of the Hyperliquid exchange totaling $195.3$ GB, covering BTC,
ETH, and SOL, with $744$ hourly files per order archive, distributed with raw book diffs and
trades. A bounded washout audit over $11$ complete hours covered $32{,}633{,}820$ events and
observed $11{,}328$ legacy order identifiers, of which $148$ were still first seen in hour
$10$. Observed-lifecycle coverage reached $99.92985\%$, but true live-book completeness at the
start of the window is unidentified. Because a valid native-FIFO test requires a certified
initial book state, we do not claim the endpoint. The decision-relevance conclusion for Q16
does not depend on it: it rests on the reconstruction diagnostic (Appendix~\ref{app:recon}),
which is computed on admitted data.

\paragraph{Provenance records.} Table~\ref{tab:prov} lists the state result identifiers, the
independent review identifiers, and the source commit for the six load-bearing packages. Each
result was accepted through an independent review before entering the programme, and the
reviews are what let us label a package as a valid negative, an inconclusive run, or a blocked
endpoint rather than asserting the status ourselves. Every package's full record, task
identifiers, protocol revisions, archived implementation branches, and artifact hashes, ships in its \texttt{provenance.json}.

\begin{table}[h]\centering
\caption{Provenance for the six load-bearing packages: state result and independent-review
identifiers. Full records, including source commits and artifact hashes, ship in each package's
\texttt{provenance.json}.}
\label{tab:prov}
\begin{tabular}{lll}
\toprule
package & results & reviews \\
\midrule
Q16 finite-queue ambiguity & R-006, R-013, R-015, R-016 & RV-004, RV-009, RV-012, RV-014 \\
Q17 forecast vs decision & R-007, R-011, R-018, R-037 & RV-004, RV-007, RV-013, RV-016 \\
Q18 depth/history/delay & R-008, R-010, R-017, R-019 & RV-004, RV-006, RV-013, RV-015 \\
Reconstruction diagnosis & R-034 & RV-032 \\
Fixed-model persistence & R-042 & RV-040 \\
BTC book-washout audit & R-044 & RV-042 \\
\bottomrule
\end{tabular}\end{table}

